\section{Discussion}
\label{sec:discussion}

The main message is simple. A measurement record contains useful structure that
the standard expectation-value pipeline does not use for training, and a
validated choice of aggregation scale recovers part of it at no quantum cost.

Adaptive shot aggregation touches only the classical side of the pipeline. The
aggregation scale $k$ sets how many shots form each training view, and the two
familiar conventions remain available as its endpoints, Raw at one shot per view
and EV at all shots per view. The scale and ridge strength are chosen from
earlier time blocks, while the held-out test block remains unseen. No circuit is
rerun and no extra shots are spent.

The method helps most in a clear regime. When a record holds many shots at each
of relatively few time steps, full averaging turns a large measurement effort
into only a handful of training examples. Partial averaging spends those same
shots differently. It produces several lightly denoised views of each target,
and the small differences among the views act as a penalty that discourages the
readout from trusting feature directions dominated by measurement noise.
Equation~\eqref{eq:loss-decomposition} makes this exact. Grouped training is
ordinary averaging plus a penalty read from the observed shot fluctuations.

Our control confirms this reading. Cov-EV applies the same type of covariance
penalty after full averaging and reaches slightly lower mean error, yet stays
close to Grouped on every task. Their similar performance supports covariance
regularization as a shared mechanism. The small gap shows that grouping is one
convenient way to obtain this regularization, not the only one, and that the
benefit arises in the classical fit rather than from additional physical
information.

The effect also transfers. It appears in the six-qubit simulation reservoir and
again in a separate ten-qubit circuit on two processors, testing transfer across
reservoir designs rather than whether more qubits help.

A practical consequence follows. A different group size can only be built from
the individual shot outcomes, so trying other scales later is possible only if
those outcomes were kept. An experiment that stores just the averaged
expectation value has already thrown away the option to regroup, while an
experiment that keeps its full shot record leaves every group size open for
free.

The hardware results have a deliberately narrow scope. Their confidence
intervals are computed from the measurement blocks of one acquisition on each
task and processor. They therefore show that the effect is stable within a
single run, not that it repeats across separate runs on different days or
calibration settings. Confirming that reproducibility is the main open question,
and the natural next tests are independent hardware campaigns, other reservoir
families, and richer readouts.

\section{Conclusion}
\label{sec:conclusion}

Turning a shot record into training data has always involved a choice, and the
field has answered it by reflex, averaging everything. We show the reflex is
beatable. Choosing the aggregation scale by validation lowers mean test error
on chaotic and real-world forecasting at the balanced budget, holds on two
superconducting processors, and requires no additional quantum cost. Because
every measurement-based reservoir produces such a record, the choice we make
explicit here is available to any of them, and the question it raises belongs
to finite-shot learning in general.
